
%%%%%%%%%%%%%%%%%%%%%%%%%%%%%%%%%%%%%%%%%%%%%%%%%%%%%%%%%%%%%
\makeatletter
\let\@oldaddcontentsline\addcontentsline
\renewcommand{\addcontentsline}[3]{}
\section*{附录}
\let\addcontentsline\@oldaddcontentsline
\makeatother

附录包含本文使用的全部代码文件、数据文件与输出图表说明，如表 \ref{tab:appendix} 所示。

\begin{table}[H]
	\centering
	\caption{附件说明表}
	\begin{tabularx}{\textwidth}{LL}
		\toprule
		文件 & 说明 \\
		\midrule
		\multicolumn{2}{l}{\textbf{代码文件}} \\
		求解/问题一/问题一\_交会定位区域直径算法.py & 楔形半平面求交、直径双算法、Thales 覆盖判据、蒙特卡洛统计 \\
		求解/问题二/问题二\_第二检测点选择策略.py & 误差传播、解析最优核对、二维布站寻优、候选区域、成本灵敏度 \\
		求解/问题三/问题三\_全向源搜索清除策略.py & 七点保证覆盖网求解、覆盖性验证、虚拟时间模型与下界、灵敏度 \\
		求解/问题四/问题四\_定向源检测清除策略.py & 半圆盘可检测域、两个判别定理验证、零漏测冗余网求解、时间代价 \\
		\midrule
		\multicolumn{2}{l}{\textbf{结果数据文件}} \\
		求解/问题一/定位区域直径结果.csv & 示范构型与覆盖失败反例的顶点坐标、直径与覆盖判定 \\
		求解/问题一/覆盖率统计表.csv & 按交会角分组的直径与覆盖率统计 \\
		求解/问题一/覆盖率对测向精度的灵敏度.csv & 覆盖率与超出量随误差上界 $\delta$ 的变化 \\
		求解/问题二/最优基线与误差结果.csv & 三种准则下的最优布站与定位误差 \\
		求解/问题二/候选区域点集.csv & 候选区域内的采样点及其极坐标描述 \\
		求解/问题二/机动成本灵敏度.csv & 最优基线随成本权重的漂移 \\
		求解/问题三/覆盖网点位表.csv & 七点保证覆盖网坐标 \\
		求解/问题三/虚拟时间下界与期望表.csv & 各源个数下的时间分项与总时间 \\
		求解/问题四/三方位检测网与对比表.csv & 七点网与 37 点冗余网的对比 \\
		求解/问题四/自适应补测代价表.csv & 不同未解频道数下的时间代价 \\
		\midrule
		\multicolumn{2}{l}{\textbf{图表文件（节选）}} \\
		问题一/figures/交会定位区域示意图.png & 定位区域、直径与直径圆的全局与局部视图 \\
		问题一/figures/直径圆覆盖率与直径分布图.png & 直径分布、覆盖率随交会角变化、精度灵敏度 \\
		问题二/figures/可行域与候选区域示意图.png & 楔形可行域、误差热力图与最优布站 \\
		问题二/figures/最优布站平面结构与成本灵敏度图.png & 平面误差结构与成本灵敏度 \\
		问题三/figures/七点保证覆盖网示意图.png & 覆盖网布局与覆盖半径曲线 \\
		问题三/figures/虚拟时间构成与灵敏度图.png & 时间分解、下界对比与侧移系数灵敏度 \\
		问题四/figures/定向源可检测域与判别定理图.png & 半圆盘可检测域与两个判别定理 \\
		问题四/figures/定向源时间惩罚与网规模权衡图.png & 时间惩罚与观测网规模权衡 \\
		\bottomrule
	\end{tabularx}
	\label{tab:appendix}
\end{table}

\noindent\textbf{运行环境：}Python 3.13，主要依赖 numpy 2.3、pandas 2.2、matplotlib 3.10、scipy 1.16。全部绘图使用 matplotlib 原生接口实现，字体采用系统中文字体回退链以保证图例中文正常显示。

\noindent\textbf{硬件环境：}Apple Silicon 平台，macOS 系统。全部脚本的最长单次运行时间不超过 60 秒。

\noindent\textbf{说明：}本题问题 3 与问题 4 的求解采用纯理论建模与算法设计口径，未接入无线电干扰源环境模拟器，因此未产生演练测试与正式测试的统计数据，正文中表 3 与表 6 的正式测试结果保留完整表头待实机测试后回填。全部算法逻辑与参数取值均已在正文中给出，可直接用于对接模拟器接口的机器狗程序实现。
